\begin{table}[htbp]
\centering
\caption{Standardized Effect Sizes}
\label{tab:sde}
\begin{tabular}{lcccccc}
\hline\hline
Outcome & $\hat{\beta}$ & SE & SD($Y$) & SDE & SE(SDE) & Classification \\
\hline
Log SFD per capita & 0.1060 & 0.0084 & 0.3756 & 0.2822 & 0.0223 & Large positive \\
Log LAT per capita & 0.5704 & 0.0746 & 5.1006 & 0.1118 & 0.0146 & Moderate positive \\
Log SFR per capita & -0.0564 & 0.0051 & 0.4793 & -0.1178 & 0.0105 & Moderate negative \\
Real balance ratio & 0.8886 & 0.1587 & 36.8131 & 0.0241 & 0.0043 & Small positive \\
\hline\hline
\end{tabular}
\begin{tablenotes}
\item \textit{Notes:} \textbf{Country:} Japan. \textbf{Research question:} Does the expiration of a temporary intergovernmental transfer grace period following municipal mergers cause fiscal distress or reveal merger efficiency gains? \textbf{Policy mechanism:} Japan's Heisei Great Mergers (1999--2010) reduced municipalities from 3,253 to 1,741; a 10+5 year Local Allocation Tax grace period inflated post-merger grants to the sum of pre-merger entitlements, with a scheduled 5-year phase-out creating a mechanical fiscal cliff at staggered dates. \textbf{Outcome definition:} Standard Fiscal Demand per capita (formula-based measure of expenditure need), implied LAT per capita (max(0, SFD $-$ SFR)), Standard Fiscal Revenue per capita (own-source tax capacity), and real balance ratio (surplus share of standard fiscal scale). \textbf{Treatment:} Binary indicator for fiscal years at or after the start of the grace period phase-out, determined mechanically by merger date. \textbf{Data:} MIC Survey on Local Public Finance and RIETI Municipality Converter, FY2011--2023, municipality-year panel of 793 cities (425 merged, 368 never-merged), 10,280 observations. \textbf{Method:} Two-way fixed effects with Sun \& Abraham (2021) interaction-weighted estimator for staggered adoption; standard errors clustered at municipality level. \textbf{Sample:} Japanese cities (\textit{shi}) that existed as post-merger entities by FY2011; towns and villages excluded due to data coverage. SDE $= \hat{\beta} / \text{SD}(Y)$ where SD($Y$) is the pre-treatment standard deviation of the outcome among merged municipalities. Classification refers to magnitude, not statistical significance: Large ($|$SDE$| > 0.15$), Moderate ($0.05$--$0.15$), Small ($0.005$--$0.05$), Null ($< 0.005$).
\end{tablenotes}
\end{table}

